%% gaol_def.tex: the packages, the page layout and the macros of the manual
%% of GAOL v5 (manual/v5/gaol.tex), which %% gaol_def.tex: the packages, the page layout and the macros of the manual
%% of GAOL v5 (manual/v5/gaol.tex), which %% gaol_def.tex: the packages, the page layout and the macros of the manual
%% of GAOL v5 (manual/v5/gaol.tex), which %% gaol_def.tex: the packages, the page layout and the macros of the manual
%% of GAOL v5 (manual/v5/gaol.tex), which \input{gaol_def} loads.
%%
%% The layout is the one of the PhD thesis of Jordan Ninin (Universite de
%% Toulouse, 2010): the class book in 12 points on A4 paper, printed on both
%% sides, the margins of 2.5 cm (top, bottom, outer) and 3.5 cm (inner), the
%% running heads of fancyhdr (the chapter on the even pages, the section on the
%% odd ones, a rule of 1 point under them, the page number between dashes at
%% the foot), the links of hyperref in blue, and the numbered environments
%% Definition, Example, Remark... of the thesis, numbered within the sections.
%% The macros of the reference (\code, \file, \pxref, the environments deffun,
%% defmethod, example, onscreen...) are those of the manual of GAOL 4 by
%% Frederic Goualard (manual/v4/manual.cls), kept under their names so that the
%% text of GAOL 4 the manual takes up, and manual/check_examples.py, which reads
%% the examples, go on working.
%%
%% Copyright (c) 2026 ENSTA, France
%%
%% Created 2026-11-05 by Jordan NININ

%%--------------------------------------------------------------------------
%% Packages
%%--------------------------------------------------------------------------
\usepackage[T1]{fontenc}
\usepackage[utf8]{inputenc}
%% Latin Modern, the outlines of Computer Modern in the T1 encoding: without
%% them pdfTeX embeds bitmaps
\usepackage{lmodern}
\usepackage{textcomp}
\usepackage[english]{babel}
\usepackage{amssymb,amsmath}
\usepackage{graphicx}
\usepackage{xcolor}
\usepackage{colortbl}
\usepackage{multirow}
\usepackage{array}
\usepackage{longtable}
\usepackage{multicol}
\usepackage{fancyvrb}
\usepackage{pifont}
\usepackage[labelsep=endash,labelfont=sc]{caption}
%% The index: \index[idx]{...}, printed by \printindex[idx]
\usepackage{index}
\newindex{idx}{idx}{ind}{Index}

%%--------------------------------------------------------------------------
%% The page
%%--------------------------------------------------------------------------
\raggedbottom
\usepackage[a4paper,top=2.5cm,bottom=2.5cm,left=3.5cm,right=2.5cm,includefoot,includehead,headheight=15pt]{geometry}

\usepackage{fancyhdr}
\pagestyle{fancy}
\fancyhead{}
\fancyfoot{}
\renewcommand{\chaptermark}[1]{\markboth{\chaptername\ \thechapter\ - \ #1}{}}
\renewcommand{\sectionmark}[1]{\markright{\thesection\ - \ #1}}
\renewcommand{\headrulewidth}{1pt}
\fancyhead[LE]{\bfseries\nouppercase{\leftmark}}
\fancyhead[RO]{\bfseries\nouppercase{\rightmark}}
\fancyfoot[CE,CO]{\bfseries--\ \thepage\ --}

%% The first page of a chapter: no running head, the page number at the foot
\fancypagestyle{plain}{%
  \fancyhead{}%
  \fancyfoot[CE,CO]{\bfseries--\ \thepage\ --}%
}
%% A page with nothing
\fancypagestyle{vide}{%
  \fancyhead{}%
  \fancyfoot{}%
  \renewcommand{\headrulewidth}{0pt}%
}
%% The page number alone, under the rule
\fancypagestyle{rien}{%
  \fancyhead{}%
  \fancyfoot[CE,CO]{\bfseries--\ \thepage\ --}%
  \renewcommand{\headrulewidth}{1pt}%
}
%% The running heads of the chapters
\fancypagestyle{YO}{%
  \fancyhead[LE]{\bfseries\nouppercase{\leftmark}}%
  \fancyhead[RO]{\bfseries\nouppercase{\rightmark}}%
  \fancyfoot[CE,CO]{\bfseries--\ \thepage\ --}%
  \renewcommand{\headrulewidth}{1pt}%
}
%% The unnumbered chapters: the same text in both running heads
\newcommand*{\fixedheads}[1]{%
  \fancyhead[LE]{\bfseries\nouppercase{#1}}%
  \fancyhead[RO]{\bfseries\nouppercase{#1}}%
  \fancyfoot[CE,CO]{\bfseries--\ \thepage\ --}%
  \renewcommand{\headrulewidth}{1pt}%
}
\fancypagestyle{INTRO}{\fixedheads{Introduction}}
\fancypagestyle{CONCLU}{\fixedheads{Conclusion}}
\fancypagestyle{BIBLIO}{\fixedheads{Bibliography}}
\fancypagestyle{INDEX}{\fixedheads{Index}}

%% The blank page before a chapter carries nothing
\makeatletter
\def\cleardoublepage{\clearpage
  \if@twoside\ifodd\c@page\else
    \hbox{}\thispagestyle{vide}\newpage\if@twocolumn\hbox{}\newpage\fi\fi\fi}
\makeatother

%% An unnumbered chapter listed in the table of contents
\newcommand{\schapter}[1]{%
  \chapter*{#1}%
  \addcontentsline{toc}{chapter}{#1}%
  \markboth{#1}{#1}%
}

%%--------------------------------------------------------------------------
%% The numbered environments of the thesis: Definition 1.2.1, Example 1.2.2...
%% numbered within the sections, with one counter for all of them
%%--------------------------------------------------------------------------
\newtheorem{Def}{Definition}[section]
\newtheorem{Lem}[Def]{Lemma}
\newtheorem{Prop}[Def]{Proposition}
\newtheorem{Th}[Def]{Theorem}
\newtheorem{Cor}[Def]{Corollary}
\newtheorem{Ex}[Def]{Example}
\newtheorem{Rem}[Def]{Remark}
\def\demo{{\em Proof. }}
\def\boite{\hfill{$\hbox{\vrule height4pt width4pt}$}}

%%--------------------------------------------------------------------------
%% Mathematics
%%--------------------------------------------------------------------------
\newcommand{\RSet}{\ensuremath{\mathbb{R}}}
\newcommand{\FSet}{\ensuremath{\mathbb{F}}}
\newcommand{\ISet}{\ensuremath{\mathbb{I}}}
\newcommand{\itv}[2]{\ensuremath{[#1,\,#2]}}
\newcommand{\roundDn}[1]{\ensuremath{\downarrow#1\downarrow}}
\newcommand{\roundUp}[1]{\ensuremath{\uparrow#1\uparrow}}
\newcommand{\roundNearest}[1]{\ensuremath{\updownarrow#1\updownarrow}}
\newcommand{\hull}[1]{\ensuremath{\mathop{\Box} #1}}
\newcommand{\leftBound}[1]{\ensuremath{\underline{#1}}}
\newcommand{\rightBound}[1]{\ensuremath{\overline{#1}}}
\newcommand*{\prevdouble}[1]{\ensuremath{{#1}^-}}
\newcommand*{\nextdouble}[1]{\ensuremath{{#1}^+}}
\newcommand{\returns}{\hbox{\ding{229}}}
\let\leq=\leqslant
\let\geq=\geqslant
\let\emptyset=\varnothing

%%--------------------------------------------------------------------------
%% Names in the text
%%--------------------------------------------------------------------------
%% Code, a name of the program
\newcommand*{\code}[1]{\textup{\texttt{#1}}}
%% A name to stand out: a value of an enumeration, a format
\newcommand*{\element}[1]{\textup{\textbf{\texttt{#1}}}}
%% A definition, in the text
\newcommand{\dfn}[1]{\textsl{#1}}
%% The name of a command or a program
\newcommand{\cmd}[1]{\textup{\texttt{#1}}}
%% The name of an option, and of a long option (-{}-name)
\newcommand{\option}[1]{\textup{\texttt{#1}}}
\newcommand{\loption}[1]{\textup{\texttt{-{}-#1}}}
%% The name of a variable or a parameter
\newcommand{\var}[1]{\textsl{#1}}
%% The name of a file
\newcommand{\file}[1]{\textup{\texttt{#1}}}
%% The name of an environment variable
\newcommand{\env}[1]{\textup{\texttt{#1}}\index[idx]{#1@\texttt{#1}}}
%% A name of IEEE 1788-2015
\newcommand*{\ieee}[1]{\textsf{#1}}
%% A non terminal and a terminal of a grammar
\newcommand*{\NT}[1]{\uppercase{\textit{#1}}}
\newcommand*{\TM}[1]{\textsf{#1}}
%% What GAOL v5 changes: a mark in the text (process.md, section 4)
\newcommand*{\newinvfive}{\unskip\nobreak\ \textcolor{blue}{\textsf{\scriptsize[GAOL~v5]}}}
%% A comment in a program, and what a program prints
\definecolor{remark}{gray}{0.40}
\newcommand{\rem}[1]{\textcolor{remark}{\textsl{#1}}}
\newcommand{\outputs}[1]{\textcolor{remark}{\textsl{\returns{} #1}}}
%% A warning, in the text: the manual writes its warnings as Remarks, this
%% command stays for the sentences that need a word
\newcommand*{\important}{\textbf{Caution.}~}

%%--------------------------------------------------------------------------
%% Cross references and citations
%%--------------------------------------------------------------------------
\newcommand*{\pxref}[1]{see Section~\hyperref[#1]{\ref{#1}}, page~\pageref{#1}}
\newcommand*{\spxref}[1]{see Section~\hyperref[#1]{[\nameref{#1}]}, page~\pageref{#1}}
\newcommand*{\xref}[1]{See Section~\hyperref[#1]{\ref{#1}}, page~\pageref{#1}}
\newcommand*{\sxref}[1]{See Section~\hyperref[#1]{[\nameref{#1}]}, page~\pageref{#1}}
\newcommand*{\exref}[1]{Section~\hyperref[#1]{\ref{#1}}, page~\pageref{#1}}
\newcommand*{\sexref}[1]{Section~\hyperref[#1]{[\nameref{#1}]}, page~\pageref{#1}}
%% The citations of the manual of GAOL 4 were printed in the margin; the
%% manual of GAOL v5 cites in the text, as the thesis does
\makeatletter
\newcommand*{\margincite}{\@ifstar{\margincite@}{\margincite@}}
\newcommand*{\margincite@}[2][]{%
  \def\@tempa{#1}%
  \ifx\@tempa\@empty\cite{#2}\else\cite[#1]{#2}\fi}
\newcommand*{\marginnocite}[1]{\nocite{#1}}
\makeatother

%%--------------------------------------------------------------------------
%% The index
%%--------------------------------------------------------------------------
\newcommand*{\findex}[1]{\index[idx]{#1@\texttt{#1}}}
\newcommand*{\cindex}[1]{\index[idx]{#1}}
\newcommand*{\idxself}{\~{}}
\newcommand*{\hyperpagebf}[1]{\textbf{\hyperpage{#1}}}
%% The main entry of a name, whose page number is in bold
\newcommand*{\defindex}[1]{\index[idx]{#1@\texttt{#1}|textbf}}

%%--------------------------------------------------------------------------
%% Programs, commands and outputs (fancyvrb): @name^argument~ is the command
%% \name{argument}, and manual/check_examples.py reads the lines @outputs^...~
%% of the examples
%%--------------------------------------------------------------------------
%% What a program prints
\DefineVerbatimEnvironment{onscreen}{Verbatim}{commandchars=\@\^\~,frame=single,label={\textsf{Output}},fontsize=\footnotesize,formatcom=\upshape}
%% An example program
\def\example{\medskip\examplevb}
\def\endexample{\endexamplevb}
\DefineVerbatimEnvironment{examplevb}{Verbatim}{fontsize=\footnotesize,commandchars=\@\^\~,frame=single,label={\textsf{Example}},numbers=left,numbersep=2pt,formatcom=\upshape}
%% A command of the shell
\DefineVerbatimEnvironment{cmdshell}{Verbatim}{commandchars=\@\^\~,fontsize=\footnotesize,formatcom=\upshape}

%%--------------------------------------------------------------------------
%% The reference: the definition of a function, a method, a type, a constant
%% or a macro, its signature on a line and its description indented under it
%%--------------------------------------------------------------------------
\makeatletter
\let\@firsttime=\@undefined
\newlength{\defwidth}
\setlength{\defwidth}{.94\textwidth}
%% A signature: the type, the name in bold, the parameters
\newcommand*{\signature}[3]{%
  \noindent{\ttfamily\raggedright\frenchspacing\hangindent=2em\hangafter=1\relax #1 \textbf{#2}#3\par}}

%% A function: #1 = return type, #2 = name, #3 = parameters
\newcommand{\deffunx}[3]{%
  \defindex{#2}%
  \ifx\@firsttime\@undefined\medskip\fi
  \gdef\@firsttime{}%
  \signature{#1}{#2}{#3}}
\newenvironment{deffun}[3]{%
  \deffunx{#1}{#2}{#3}%
  \begin{flushright}%
  \begin{minipage}{\defwidth}%
}{%
  \end{minipage}%
  \end{flushright}%
  \medskip
  \global\let\@firsttime=\@undefined
}

%% A function template: #1 = template, #2 = return type, #3 = name,
%% #4 = parameters
\newcommand{\defTfunx}[4]{%
  \defindex{#3}%
  \ifx\@firsttime\@undefined\medskip\fi
  \gdef\@firsttime{}%
  \noindent{\ttfamily\raggedright\frenchspacing #1\par}
  \signature{#2}{#3}{#4}}
\newenvironment{defTfun}[4]{%
  \defTfunx{#1}{#2}{#3}{#4}%
  \begin{flushright}%
  \begin{minipage}{\defwidth}%
}{%
  \end{minipage}%
  \end{flushright}%
  \medskip
  \global\let\@firsttime=\@undefined
}

%% A method: #1 = class, #2 = return type, #3 = name, #4 = parameters
\newcommand{\defmethodx}[4]{%
  \defindex{#3}%
  \ifx\@firsttime\@undefined\medskip\fi
  \gdef\@firsttime{}%
  \noindent{\ttfamily\raggedright\frenchspacing\hangindent=2em\hangafter=1\relax\if#2\@empty\else #2 \fi#1::\textbf{#3}#4\par}}
\newenvironment{defmethod}[4]{%
  \defmethodx{#1}{#2}{#3}{#4}%
  \begin{flushright}%
  \begin{minipage}{\defwidth}%
}{%
  \end{minipage}%
  \end{flushright}%
  \medskip
  \global\let\@firsttime=\@undefined
}

%% A static method: #1 = class, #2 = return type, #3 = name, #4 = parameters
\newcommand{\defstaticmethodx}[4]{%
  \defindex{#3}%
  \ifx\@firsttime\@undefined\medskip\fi
  \gdef\@firsttime{}%
  \noindent{\ttfamily\frenchspacing #2 #1::\textbf{#3}#4\hfill\textrm{\textit{static}}}\par}
\newenvironment{defstaticmethod}[4]{%
  \defstaticmethodx{#1}{#2}{#3}{#4}%
  \begin{flushright}%
  \begin{minipage}{\defwidth}%
}{%
  \end{minipage}%
  \end{flushright}%
  \medskip
  \global\let\@firsttime=\@undefined
}

%% A type: #1 = kind (struct, class, enum), #2 = name
\newenvironment{deftype}[2]{%
  \defindex{#2}%
  \medskip
  \noindent{\ttfamily\textbf{#2}\hfill\textrm{\textit{#1}}}\par
  \begin{flushright}%
  \begin{minipage}{\defwidth}%
}{%
  \end{minipage}%
  \end{flushright}%
  \medskip
}

%% A constant: #1 = type, #2 = name
\newenvironment{defconst}[2]{%
  \defindex{#2}%
  \medskip
  \noindent{\ttfamily\frenchspacing #1 \textbf{#2}\hfill\textrm{\textit{const}}}\par
  \begin{flushright}%
  \begin{minipage}{\defwidth}%
}{%
  \end{minipage}%
  \end{flushright}%
  \medskip
}

%% A macro without parameter: #1 = name
\newenvironment{defmacro}[1]{%
  \defindex{#1}%
  \medskip
  \noindent{\ttfamily\textbf{#1}\hfill\textrm{\textit{macro}}}\par
  \begin{flushright}%
  \begin{minipage}{\defwidth}%
}{%
  \end{minipage}%
  \end{flushright}%
  \medskip
}

%% A macro with parameters: #1 = name, #2 = parameters
\newenvironment{defmacroA}[2]{%
  \defindex{#1}%
  \medskip
  \noindent{\ttfamily\textbf{#1}#2\hfill\textrm{\textit{macro}}}\par
  \begin{flushright}%
  \begin{minipage}{\defwidth}%
}{%
  \end{minipage}%
  \end{flushright}%
  \medskip
}

%% The mathematical definition of an operation, one \op per line
\newenvironment{operation}{%
  \vspace*{-5pt}%
  \def\op##1{\noindent{\small\ensuremath{\blacksquare}}\ ##1\par}%
}{%
  \par\smallskip
}
\makeatother

%% The page of the copyright notice
\newenvironment{copyrighttxt}{%
  \parindent0pt%
  \newpage
  \thispagestyle{empty}%
  \null\vfill
}{%
  \newpage
}

%% The numbers of the table of contents and of the lists are wider than in the
%% thesis (3.10.10, 12.13)
\makeatletter
\renewcommand*{\l@section}{\@dottedtocline{1}{1.5em}{2.8em}}
\renewcommand*{\l@subsection}{\@dottedtocline{2}{4.3em}{3.6em}}
\renewcommand*{\l@table}{\@dottedtocline{1}{1.5em}{3.0em}}
\renewcommand*{\l@figure}{\@dottedtocline{1}{1.5em}{3.0em}}
\makeatother

%% A column of a table whose text is set ragged right: L{width}
\newcolumntype{L}[1]{>{\raggedright\arraybackslash}p{#1}}
%% The names of functions and options, in typewriter, are not hyphenated: let
%% the lines holding them stretch rather than overflow into the margin
\setlength{\emergencystretch}{3em}
%% A box that overflows the margin by less than 5 pt, which is 2.5 cm wide, is
%% not worth a warning
\hfuzz=5pt

%%--------------------------------------------------------------------------
%% hyperref comes last
%%--------------------------------------------------------------------------
\usepackage{hyperref}
%% The names of macros and functions, in typewriter, are long and are not
%% hyphenated: a line may break after an underscore (GAOL_PRESERVE_ROUNDING),
%% without a hyphen. In the strings of the PDF the underscore stays itself.
\makeatletter
\let\gaol@underscore\_
\renewcommand*{\_}{\texorpdfstring{\gaol@underscore\hspace{0pt}}{\gaol@underscore}}
\makeatother
\hypersetup{
    unicode      = false,
    pdftoolbar   = true,
    pdfmenubar   = true,
    pdffitwindow = true,
    pdftitle     = {GAOL v5: NOT Just Another Interval Arithmetic Library},
    pdfauthor    = {Jordan Ninin, Frederic Goualard},
    pdfsubject   = {The manual of GAOL v5},
    pdfnewwindow = true,
    pdfkeywords  = {interval arithmetic, IEEE 1788-2015, CORE-MATH,
                    constraint programming, C++},
    colorlinks   = true,
    linkcolor    = blue,
    citecolor    = blue,
    filecolor    = blue,
    urlcolor     = blue,
    breaklinks   = true
}
 loads.
%%
%% The layout is the one of the PhD thesis of Jordan Ninin (Universite de
%% Toulouse, 2010): the class book in 12 points on A4 paper, printed on both
%% sides, the margins of 2.5 cm (top, bottom, outer) and 3.5 cm (inner), the
%% running heads of fancyhdr (the chapter on the even pages, the section on the
%% odd ones, a rule of 1 point under them, the page number between dashes at
%% the foot), the links of hyperref in blue, and the numbered environments
%% Definition, Example, Remark... of the thesis, numbered within the sections.
%% The macros of the reference (\code, \file, \pxref, the environments deffun,
%% defmethod, example, onscreen...) are those of the manual of GAOL 4 by
%% Frederic Goualard (manual/v4/manual.cls), kept under their names so that the
%% text of GAOL 4 the manual takes up, and manual/check_examples.py, which reads
%% the examples, go on working.
%%
%% Copyright (c) 2026 ENSTA, France
%%
%% Created 2026-11-05 by Jordan NININ

%%--------------------------------------------------------------------------
%% Packages
%%--------------------------------------------------------------------------
\usepackage[T1]{fontenc}
\usepackage[utf8]{inputenc}
%% Latin Modern, the outlines of Computer Modern in the T1 encoding: without
%% them pdfTeX embeds bitmaps
\usepackage{lmodern}
\usepackage{textcomp}
\usepackage[english]{babel}
\usepackage{amssymb,amsmath}
\usepackage{graphicx}
\usepackage{xcolor}
\usepackage{colortbl}
\usepackage{multirow}
\usepackage{array}
\usepackage{longtable}
\usepackage{multicol}
\usepackage{fancyvrb}
\usepackage{pifont}
\usepackage[labelsep=endash,labelfont=sc]{caption}
%% The index: \index[idx]{...}, printed by \printindex[idx]
\usepackage{index}
\newindex{idx}{idx}{ind}{Index}

%%--------------------------------------------------------------------------
%% The page
%%--------------------------------------------------------------------------
\raggedbottom
\usepackage[a4paper,top=2.5cm,bottom=2.5cm,left=3.5cm,right=2.5cm,includefoot,includehead,headheight=15pt]{geometry}

\usepackage{fancyhdr}
\pagestyle{fancy}
\fancyhead{}
\fancyfoot{}
\renewcommand{\chaptermark}[1]{\markboth{\chaptername\ \thechapter\ - \ #1}{}}
\renewcommand{\sectionmark}[1]{\markright{\thesection\ - \ #1}}
\renewcommand{\headrulewidth}{1pt}
\fancyhead[LE]{\bfseries\nouppercase{\leftmark}}
\fancyhead[RO]{\bfseries\nouppercase{\rightmark}}
\fancyfoot[CE,CO]{\bfseries--\ \thepage\ --}

%% The first page of a chapter: no running head, the page number at the foot
\fancypagestyle{plain}{%
  \fancyhead{}%
  \fancyfoot[CE,CO]{\bfseries--\ \thepage\ --}%
}
%% A page with nothing
\fancypagestyle{vide}{%
  \fancyhead{}%
  \fancyfoot{}%
  \renewcommand{\headrulewidth}{0pt}%
}
%% The page number alone, under the rule
\fancypagestyle{rien}{%
  \fancyhead{}%
  \fancyfoot[CE,CO]{\bfseries--\ \thepage\ --}%
  \renewcommand{\headrulewidth}{1pt}%
}
%% The running heads of the chapters
\fancypagestyle{YO}{%
  \fancyhead[LE]{\bfseries\nouppercase{\leftmark}}%
  \fancyhead[RO]{\bfseries\nouppercase{\rightmark}}%
  \fancyfoot[CE,CO]{\bfseries--\ \thepage\ --}%
  \renewcommand{\headrulewidth}{1pt}%
}
%% The unnumbered chapters: the same text in both running heads
\newcommand*{\fixedheads}[1]{%
  \fancyhead[LE]{\bfseries\nouppercase{#1}}%
  \fancyhead[RO]{\bfseries\nouppercase{#1}}%
  \fancyfoot[CE,CO]{\bfseries--\ \thepage\ --}%
  \renewcommand{\headrulewidth}{1pt}%
}
\fancypagestyle{INTRO}{\fixedheads{Introduction}}
\fancypagestyle{CONCLU}{\fixedheads{Conclusion}}
\fancypagestyle{BIBLIO}{\fixedheads{Bibliography}}
\fancypagestyle{INDEX}{\fixedheads{Index}}

%% The blank page before a chapter carries nothing
\makeatletter
\def\cleardoublepage{\clearpage
  \if@twoside\ifodd\c@page\else
    \hbox{}\thispagestyle{vide}\newpage\if@twocolumn\hbox{}\newpage\fi\fi\fi}
\makeatother

%% An unnumbered chapter listed in the table of contents
\newcommand{\schapter}[1]{%
  \chapter*{#1}%
  \addcontentsline{toc}{chapter}{#1}%
  \markboth{#1}{#1}%
}

%%--------------------------------------------------------------------------
%% The numbered environments of the thesis: Definition 1.2.1, Example 1.2.2...
%% numbered within the sections, with one counter for all of them
%%--------------------------------------------------------------------------
\newtheorem{Def}{Definition}[section]
\newtheorem{Lem}[Def]{Lemma}
\newtheorem{Prop}[Def]{Proposition}
\newtheorem{Th}[Def]{Theorem}
\newtheorem{Cor}[Def]{Corollary}
\newtheorem{Ex}[Def]{Example}
\newtheorem{Rem}[Def]{Remark}
\def\demo{{\em Proof. }}
\def\boite{\hfill{$\hbox{\vrule height4pt width4pt}$}}

%%--------------------------------------------------------------------------
%% Mathematics
%%--------------------------------------------------------------------------
\newcommand{\RSet}{\ensuremath{\mathbb{R}}}
\newcommand{\FSet}{\ensuremath{\mathbb{F}}}
\newcommand{\ISet}{\ensuremath{\mathbb{I}}}
\newcommand{\itv}[2]{\ensuremath{[#1,\,#2]}}
\newcommand{\roundDn}[1]{\ensuremath{\downarrow#1\downarrow}}
\newcommand{\roundUp}[1]{\ensuremath{\uparrow#1\uparrow}}
\newcommand{\roundNearest}[1]{\ensuremath{\updownarrow#1\updownarrow}}
\newcommand{\hull}[1]{\ensuremath{\mathop{\Box} #1}}
\newcommand{\leftBound}[1]{\ensuremath{\underline{#1}}}
\newcommand{\rightBound}[1]{\ensuremath{\overline{#1}}}
\newcommand*{\prevdouble}[1]{\ensuremath{{#1}^-}}
\newcommand*{\nextdouble}[1]{\ensuremath{{#1}^+}}
\newcommand{\returns}{\hbox{\ding{229}}}
\let\leq=\leqslant
\let\geq=\geqslant
\let\emptyset=\varnothing

%%--------------------------------------------------------------------------
%% Names in the text
%%--------------------------------------------------------------------------
%% Code, a name of the program
\newcommand*{\code}[1]{\textup{\texttt{#1}}}
%% A name to stand out: a value of an enumeration, a format
\newcommand*{\element}[1]{\textup{\textbf{\texttt{#1}}}}
%% A definition, in the text
\newcommand{\dfn}[1]{\textsl{#1}}
%% The name of a command or a program
\newcommand{\cmd}[1]{\textup{\texttt{#1}}}
%% The name of an option, and of a long option (-{}-name)
\newcommand{\option}[1]{\textup{\texttt{#1}}}
\newcommand{\loption}[1]{\textup{\texttt{-{}-#1}}}
%% The name of a variable or a parameter
\newcommand{\var}[1]{\textsl{#1}}
%% The name of a file
\newcommand{\file}[1]{\textup{\texttt{#1}}}
%% The name of an environment variable
\newcommand{\env}[1]{\textup{\texttt{#1}}\index[idx]{#1@\texttt{#1}}}
%% A name of IEEE 1788-2015
\newcommand*{\ieee}[1]{\textsf{#1}}
%% A non terminal and a terminal of a grammar
\newcommand*{\NT}[1]{\uppercase{\textit{#1}}}
\newcommand*{\TM}[1]{\textsf{#1}}
%% What GAOL v5 changes: a mark in the text (process.md, section 4)
\newcommand*{\newinvfive}{\unskip\nobreak\ \textcolor{blue}{\textsf{\scriptsize[GAOL~v5]}}}
%% A comment in a program, and what a program prints
\definecolor{remark}{gray}{0.40}
\newcommand{\rem}[1]{\textcolor{remark}{\textsl{#1}}}
\newcommand{\outputs}[1]{\textcolor{remark}{\textsl{\returns{} #1}}}
%% A warning, in the text: the manual writes its warnings as Remarks, this
%% command stays for the sentences that need a word
\newcommand*{\important}{\textbf{Caution.}~}

%%--------------------------------------------------------------------------
%% Cross references and citations
%%--------------------------------------------------------------------------
\newcommand*{\pxref}[1]{see Section~\hyperref[#1]{\ref{#1}}, page~\pageref{#1}}
\newcommand*{\spxref}[1]{see Section~\hyperref[#1]{[\nameref{#1}]}, page~\pageref{#1}}
\newcommand*{\xref}[1]{See Section~\hyperref[#1]{\ref{#1}}, page~\pageref{#1}}
\newcommand*{\sxref}[1]{See Section~\hyperref[#1]{[\nameref{#1}]}, page~\pageref{#1}}
\newcommand*{\exref}[1]{Section~\hyperref[#1]{\ref{#1}}, page~\pageref{#1}}
\newcommand*{\sexref}[1]{Section~\hyperref[#1]{[\nameref{#1}]}, page~\pageref{#1}}
%% The citations of the manual of GAOL 4 were printed in the margin; the
%% manual of GAOL v5 cites in the text, as the thesis does
\makeatletter
\newcommand*{\margincite}{\@ifstar{\margincite@}{\margincite@}}
\newcommand*{\margincite@}[2][]{%
  \def\@tempa{#1}%
  \ifx\@tempa\@empty\cite{#2}\else\cite[#1]{#2}\fi}
\newcommand*{\marginnocite}[1]{\nocite{#1}}
\makeatother

%%--------------------------------------------------------------------------
%% The index
%%--------------------------------------------------------------------------
\newcommand*{\findex}[1]{\index[idx]{#1@\texttt{#1}}}
\newcommand*{\cindex}[1]{\index[idx]{#1}}
\newcommand*{\idxself}{\~{}}
\newcommand*{\hyperpagebf}[1]{\textbf{\hyperpage{#1}}}
%% The main entry of a name, whose page number is in bold
\newcommand*{\defindex}[1]{\index[idx]{#1@\texttt{#1}|textbf}}

%%--------------------------------------------------------------------------
%% Programs, commands and outputs (fancyvrb): @name^argument~ is the command
%% \name{argument}, and manual/check_examples.py reads the lines @outputs^...~
%% of the examples
%%--------------------------------------------------------------------------
%% What a program prints
\DefineVerbatimEnvironment{onscreen}{Verbatim}{commandchars=\@\^\~,frame=single,label={\textsf{Output}},fontsize=\footnotesize,formatcom=\upshape}
%% An example program
\def\example{\medskip\examplevb}
\def\endexample{\endexamplevb}
\DefineVerbatimEnvironment{examplevb}{Verbatim}{fontsize=\footnotesize,commandchars=\@\^\~,frame=single,label={\textsf{Example}},numbers=left,numbersep=2pt,formatcom=\upshape}
%% A command of the shell
\DefineVerbatimEnvironment{cmdshell}{Verbatim}{commandchars=\@\^\~,fontsize=\footnotesize,formatcom=\upshape}

%%--------------------------------------------------------------------------
%% The reference: the definition of a function, a method, a type, a constant
%% or a macro, its signature on a line and its description indented under it
%%--------------------------------------------------------------------------
\makeatletter
\let\@firsttime=\@undefined
\newlength{\defwidth}
\setlength{\defwidth}{.94\textwidth}
%% A signature: the type, the name in bold, the parameters
\newcommand*{\signature}[3]{%
  \noindent{\ttfamily\raggedright\frenchspacing\hangindent=2em\hangafter=1\relax #1 \textbf{#2}#3\par}}

%% A function: #1 = return type, #2 = name, #3 = parameters
\newcommand{\deffunx}[3]{%
  \defindex{#2}%
  \ifx\@firsttime\@undefined\medskip\fi
  \gdef\@firsttime{}%
  \signature{#1}{#2}{#3}}
\newenvironment{deffun}[3]{%
  \deffunx{#1}{#2}{#3}%
  \begin{flushright}%
  \begin{minipage}{\defwidth}%
}{%
  \end{minipage}%
  \end{flushright}%
  \medskip
  \global\let\@firsttime=\@undefined
}

%% A function template: #1 = template, #2 = return type, #3 = name,
%% #4 = parameters
\newcommand{\defTfunx}[4]{%
  \defindex{#3}%
  \ifx\@firsttime\@undefined\medskip\fi
  \gdef\@firsttime{}%
  \noindent{\ttfamily\raggedright\frenchspacing #1\par}
  \signature{#2}{#3}{#4}}
\newenvironment{defTfun}[4]{%
  \defTfunx{#1}{#2}{#3}{#4}%
  \begin{flushright}%
  \begin{minipage}{\defwidth}%
}{%
  \end{minipage}%
  \end{flushright}%
  \medskip
  \global\let\@firsttime=\@undefined
}

%% A method: #1 = class, #2 = return type, #3 = name, #4 = parameters
\newcommand{\defmethodx}[4]{%
  \defindex{#3}%
  \ifx\@firsttime\@undefined\medskip\fi
  \gdef\@firsttime{}%
  \noindent{\ttfamily\raggedright\frenchspacing\hangindent=2em\hangafter=1\relax\if#2\@empty\else #2 \fi#1::\textbf{#3}#4\par}}
\newenvironment{defmethod}[4]{%
  \defmethodx{#1}{#2}{#3}{#4}%
  \begin{flushright}%
  \begin{minipage}{\defwidth}%
}{%
  \end{minipage}%
  \end{flushright}%
  \medskip
  \global\let\@firsttime=\@undefined
}

%% A static method: #1 = class, #2 = return type, #3 = name, #4 = parameters
\newcommand{\defstaticmethodx}[4]{%
  \defindex{#3}%
  \ifx\@firsttime\@undefined\medskip\fi
  \gdef\@firsttime{}%
  \noindent{\ttfamily\frenchspacing #2 #1::\textbf{#3}#4\hfill\textrm{\textit{static}}}\par}
\newenvironment{defstaticmethod}[4]{%
  \defstaticmethodx{#1}{#2}{#3}{#4}%
  \begin{flushright}%
  \begin{minipage}{\defwidth}%
}{%
  \end{minipage}%
  \end{flushright}%
  \medskip
  \global\let\@firsttime=\@undefined
}

%% A type: #1 = kind (struct, class, enum), #2 = name
\newenvironment{deftype}[2]{%
  \defindex{#2}%
  \medskip
  \noindent{\ttfamily\textbf{#2}\hfill\textrm{\textit{#1}}}\par
  \begin{flushright}%
  \begin{minipage}{\defwidth}%
}{%
  \end{minipage}%
  \end{flushright}%
  \medskip
}

%% A constant: #1 = type, #2 = name
\newenvironment{defconst}[2]{%
  \defindex{#2}%
  \medskip
  \noindent{\ttfamily\frenchspacing #1 \textbf{#2}\hfill\textrm{\textit{const}}}\par
  \begin{flushright}%
  \begin{minipage}{\defwidth}%
}{%
  \end{minipage}%
  \end{flushright}%
  \medskip
}

%% A macro without parameter: #1 = name
\newenvironment{defmacro}[1]{%
  \defindex{#1}%
  \medskip
  \noindent{\ttfamily\textbf{#1}\hfill\textrm{\textit{macro}}}\par
  \begin{flushright}%
  \begin{minipage}{\defwidth}%
}{%
  \end{minipage}%
  \end{flushright}%
  \medskip
}

%% A macro with parameters: #1 = name, #2 = parameters
\newenvironment{defmacroA}[2]{%
  \defindex{#1}%
  \medskip
  \noindent{\ttfamily\textbf{#1}#2\hfill\textrm{\textit{macro}}}\par
  \begin{flushright}%
  \begin{minipage}{\defwidth}%
}{%
  \end{minipage}%
  \end{flushright}%
  \medskip
}

%% The mathematical definition of an operation, one \op per line
\newenvironment{operation}{%
  \vspace*{-5pt}%
  \def\op##1{\noindent{\small\ensuremath{\blacksquare}}\ ##1\par}%
}{%
  \par\smallskip
}
\makeatother

%% The page of the copyright notice
\newenvironment{copyrighttxt}{%
  \parindent0pt%
  \newpage
  \thispagestyle{empty}%
  \null\vfill
}{%
  \newpage
}

%% The numbers of the table of contents and of the lists are wider than in the
%% thesis (3.10.10, 12.13)
\makeatletter
\renewcommand*{\l@section}{\@dottedtocline{1}{1.5em}{2.8em}}
\renewcommand*{\l@subsection}{\@dottedtocline{2}{4.3em}{3.6em}}
\renewcommand*{\l@table}{\@dottedtocline{1}{1.5em}{3.0em}}
\renewcommand*{\l@figure}{\@dottedtocline{1}{1.5em}{3.0em}}
\makeatother

%% A column of a table whose text is set ragged right: L{width}
\newcolumntype{L}[1]{>{\raggedright\arraybackslash}p{#1}}
%% The names of functions and options, in typewriter, are not hyphenated: let
%% the lines holding them stretch rather than overflow into the margin
\setlength{\emergencystretch}{3em}
%% A box that overflows the margin by less than 5 pt, which is 2.5 cm wide, is
%% not worth a warning
\hfuzz=5pt

%%--------------------------------------------------------------------------
%% hyperref comes last
%%--------------------------------------------------------------------------
\usepackage{hyperref}
%% The names of macros and functions, in typewriter, are long and are not
%% hyphenated: a line may break after an underscore (GAOL_PRESERVE_ROUNDING),
%% without a hyphen. In the strings of the PDF the underscore stays itself.
\makeatletter
\let\gaol@underscore\_
\renewcommand*{\_}{\texorpdfstring{\gaol@underscore\hspace{0pt}}{\gaol@underscore}}
\makeatother
\hypersetup{
    unicode      = false,
    pdftoolbar   = true,
    pdfmenubar   = true,
    pdffitwindow = true,
    pdftitle     = {GAOL v5: NOT Just Another Interval Arithmetic Library},
    pdfauthor    = {Jordan Ninin, Frederic Goualard},
    pdfsubject   = {The manual of GAOL v5},
    pdfnewwindow = true,
    pdfkeywords  = {interval arithmetic, IEEE 1788-2015, CORE-MATH,
                    constraint programming, C++},
    colorlinks   = true,
    linkcolor    = blue,
    citecolor    = blue,
    filecolor    = blue,
    urlcolor     = blue,
    breaklinks   = true
}
 loads.
%%
%% The layout is the one of the PhD thesis of Jordan Ninin (Universite de
%% Toulouse, 2010): the class book in 12 points on A4 paper, printed on both
%% sides, the margins of 2.5 cm (top, bottom, outer) and 3.5 cm (inner), the
%% running heads of fancyhdr (the chapter on the even pages, the section on the
%% odd ones, a rule of 1 point under them, the page number between dashes at
%% the foot), the links of hyperref in blue, and the numbered environments
%% Definition, Example, Remark... of the thesis, numbered within the sections.
%% The macros of the reference (\code, \file, \pxref, the environments deffun,
%% defmethod, example, onscreen...) are those of the manual of GAOL 4 by
%% Frederic Goualard (manual/v4/manual.cls), kept under their names so that the
%% text of GAOL 4 the manual takes up, and manual/check_examples.py, which reads
%% the examples, go on working.
%%
%% Copyright (c) 2026 ENSTA, France
%%
%% Created 2026-11-05 by Jordan NININ

%%--------------------------------------------------------------------------
%% Packages
%%--------------------------------------------------------------------------
\usepackage[T1]{fontenc}
\usepackage[utf8]{inputenc}
%% Latin Modern, the outlines of Computer Modern in the T1 encoding: without
%% them pdfTeX embeds bitmaps
\usepackage{lmodern}
\usepackage{textcomp}
\usepackage[english]{babel}
\usepackage{amssymb,amsmath}
\usepackage{graphicx}
\usepackage{xcolor}
\usepackage{colortbl}
\usepackage{multirow}
\usepackage{array}
\usepackage{longtable}
\usepackage{multicol}
\usepackage{fancyvrb}
\usepackage{pifont}
\usepackage[labelsep=endash,labelfont=sc]{caption}
%% The index: \index[idx]{...}, printed by \printindex[idx]
\usepackage{index}
\newindex{idx}{idx}{ind}{Index}

%%--------------------------------------------------------------------------
%% The page
%%--------------------------------------------------------------------------
\raggedbottom
\usepackage[a4paper,top=2.5cm,bottom=2.5cm,left=3.5cm,right=2.5cm,includefoot,includehead,headheight=15pt]{geometry}

\usepackage{fancyhdr}
\pagestyle{fancy}
\fancyhead{}
\fancyfoot{}
\renewcommand{\chaptermark}[1]{\markboth{\chaptername\ \thechapter\ - \ #1}{}}
\renewcommand{\sectionmark}[1]{\markright{\thesection\ - \ #1}}
\renewcommand{\headrulewidth}{1pt}
\fancyhead[LE]{\bfseries\nouppercase{\leftmark}}
\fancyhead[RO]{\bfseries\nouppercase{\rightmark}}
\fancyfoot[CE,CO]{\bfseries--\ \thepage\ --}

%% The first page of a chapter: no running head, the page number at the foot
\fancypagestyle{plain}{%
  \fancyhead{}%
  \fancyfoot[CE,CO]{\bfseries--\ \thepage\ --}%
}
%% A page with nothing
\fancypagestyle{vide}{%
  \fancyhead{}%
  \fancyfoot{}%
  \renewcommand{\headrulewidth}{0pt}%
}
%% The page number alone, under the rule
\fancypagestyle{rien}{%
  \fancyhead{}%
  \fancyfoot[CE,CO]{\bfseries--\ \thepage\ --}%
  \renewcommand{\headrulewidth}{1pt}%
}
%% The running heads of the chapters
\fancypagestyle{YO}{%
  \fancyhead[LE]{\bfseries\nouppercase{\leftmark}}%
  \fancyhead[RO]{\bfseries\nouppercase{\rightmark}}%
  \fancyfoot[CE,CO]{\bfseries--\ \thepage\ --}%
  \renewcommand{\headrulewidth}{1pt}%
}
%% The unnumbered chapters: the same text in both running heads
\newcommand*{\fixedheads}[1]{%
  \fancyhead[LE]{\bfseries\nouppercase{#1}}%
  \fancyhead[RO]{\bfseries\nouppercase{#1}}%
  \fancyfoot[CE,CO]{\bfseries--\ \thepage\ --}%
  \renewcommand{\headrulewidth}{1pt}%
}
\fancypagestyle{INTRO}{\fixedheads{Introduction}}
\fancypagestyle{CONCLU}{\fixedheads{Conclusion}}
\fancypagestyle{BIBLIO}{\fixedheads{Bibliography}}
\fancypagestyle{INDEX}{\fixedheads{Index}}

%% The blank page before a chapter carries nothing
\makeatletter
\def\cleardoublepage{\clearpage
  \if@twoside\ifodd\c@page\else
    \hbox{}\thispagestyle{vide}\newpage\if@twocolumn\hbox{}\newpage\fi\fi\fi}
\makeatother

%% An unnumbered chapter listed in the table of contents
\newcommand{\schapter}[1]{%
  \chapter*{#1}%
  \addcontentsline{toc}{chapter}{#1}%
  \markboth{#1}{#1}%
}

%%--------------------------------------------------------------------------
%% The numbered environments of the thesis: Definition 1.2.1, Example 1.2.2...
%% numbered within the sections, with one counter for all of them
%%--------------------------------------------------------------------------
\newtheorem{Def}{Definition}[section]
\newtheorem{Lem}[Def]{Lemma}
\newtheorem{Prop}[Def]{Proposition}
\newtheorem{Th}[Def]{Theorem}
\newtheorem{Cor}[Def]{Corollary}
\newtheorem{Ex}[Def]{Example}
\newtheorem{Rem}[Def]{Remark}
\def\demo{{\em Proof. }}
\def\boite{\hfill{$\hbox{\vrule height4pt width4pt}$}}

%%--------------------------------------------------------------------------
%% Mathematics
%%--------------------------------------------------------------------------
\newcommand{\RSet}{\ensuremath{\mathbb{R}}}
\newcommand{\FSet}{\ensuremath{\mathbb{F}}}
\newcommand{\ISet}{\ensuremath{\mathbb{I}}}
\newcommand{\itv}[2]{\ensuremath{[#1,\,#2]}}
\newcommand{\roundDn}[1]{\ensuremath{\downarrow#1\downarrow}}
\newcommand{\roundUp}[1]{\ensuremath{\uparrow#1\uparrow}}
\newcommand{\roundNearest}[1]{\ensuremath{\updownarrow#1\updownarrow}}
\newcommand{\hull}[1]{\ensuremath{\mathop{\Box} #1}}
\newcommand{\leftBound}[1]{\ensuremath{\underline{#1}}}
\newcommand{\rightBound}[1]{\ensuremath{\overline{#1}}}
\newcommand*{\prevdouble}[1]{\ensuremath{{#1}^-}}
\newcommand*{\nextdouble}[1]{\ensuremath{{#1}^+}}
\newcommand{\returns}{\hbox{\ding{229}}}
\let\leq=\leqslant
\let\geq=\geqslant
\let\emptyset=\varnothing

%%--------------------------------------------------------------------------
%% Names in the text
%%--------------------------------------------------------------------------
%% Code, a name of the program
\newcommand*{\code}[1]{\textup{\texttt{#1}}}
%% A name to stand out: a value of an enumeration, a format
\newcommand*{\element}[1]{\textup{\textbf{\texttt{#1}}}}
%% A definition, in the text
\newcommand{\dfn}[1]{\textsl{#1}}
%% The name of a command or a program
\newcommand{\cmd}[1]{\textup{\texttt{#1}}}
%% The name of an option, and of a long option (-{}-name)
\newcommand{\option}[1]{\textup{\texttt{#1}}}
\newcommand{\loption}[1]{\textup{\texttt{-{}-#1}}}
%% The name of a variable or a parameter
\newcommand{\var}[1]{\textsl{#1}}
%% The name of a file
\newcommand{\file}[1]{\textup{\texttt{#1}}}
%% The name of an environment variable
\newcommand{\env}[1]{\textup{\texttt{#1}}\index[idx]{#1@\texttt{#1}}}
%% A name of IEEE 1788-2015
\newcommand*{\ieee}[1]{\textsf{#1}}
%% A non terminal and a terminal of a grammar
\newcommand*{\NT}[1]{\uppercase{\textit{#1}}}
\newcommand*{\TM}[1]{\textsf{#1}}
%% What GAOL v5 changes: a mark in the text (process.md, section 4)
\newcommand*{\newinvfive}{\unskip\nobreak\ \textcolor{blue}{\textsf{\scriptsize[GAOL~v5]}}}
%% A comment in a program, and what a program prints
\definecolor{remark}{gray}{0.40}
\newcommand{\rem}[1]{\textcolor{remark}{\textsl{#1}}}
\newcommand{\outputs}[1]{\textcolor{remark}{\textsl{\returns{} #1}}}
%% A warning, in the text: the manual writes its warnings as Remarks, this
%% command stays for the sentences that need a word
\newcommand*{\important}{\textbf{Caution.}~}

%%--------------------------------------------------------------------------
%% Cross references and citations
%%--------------------------------------------------------------------------
\newcommand*{\pxref}[1]{see Section~\hyperref[#1]{\ref{#1}}, page~\pageref{#1}}
\newcommand*{\spxref}[1]{see Section~\hyperref[#1]{[\nameref{#1}]}, page~\pageref{#1}}
\newcommand*{\xref}[1]{See Section~\hyperref[#1]{\ref{#1}}, page~\pageref{#1}}
\newcommand*{\sxref}[1]{See Section~\hyperref[#1]{[\nameref{#1}]}, page~\pageref{#1}}
\newcommand*{\exref}[1]{Section~\hyperref[#1]{\ref{#1}}, page~\pageref{#1}}
\newcommand*{\sexref}[1]{Section~\hyperref[#1]{[\nameref{#1}]}, page~\pageref{#1}}
%% The citations of the manual of GAOL 4 were printed in the margin; the
%% manual of GAOL v5 cites in the text, as the thesis does
\makeatletter
\newcommand*{\margincite}{\@ifstar{\margincite@}{\margincite@}}
\newcommand*{\margincite@}[2][]{%
  \def\@tempa{#1}%
  \ifx\@tempa\@empty\cite{#2}\else\cite[#1]{#2}\fi}
\newcommand*{\marginnocite}[1]{\nocite{#1}}
\makeatother

%%--------------------------------------------------------------------------
%% The index
%%--------------------------------------------------------------------------
\newcommand*{\findex}[1]{\index[idx]{#1@\texttt{#1}}}
\newcommand*{\cindex}[1]{\index[idx]{#1}}
\newcommand*{\idxself}{\~{}}
\newcommand*{\hyperpagebf}[1]{\textbf{\hyperpage{#1}}}
%% The main entry of a name, whose page number is in bold
\newcommand*{\defindex}[1]{\index[idx]{#1@\texttt{#1}|textbf}}

%%--------------------------------------------------------------------------
%% Programs, commands and outputs (fancyvrb): @name^argument~ is the command
%% \name{argument}, and manual/check_examples.py reads the lines @outputs^...~
%% of the examples
%%--------------------------------------------------------------------------
%% What a program prints
\DefineVerbatimEnvironment{onscreen}{Verbatim}{commandchars=\@\^\~,frame=single,label={\textsf{Output}},fontsize=\footnotesize,formatcom=\upshape}
%% An example program
\def\example{\medskip\examplevb}
\def\endexample{\endexamplevb}
\DefineVerbatimEnvironment{examplevb}{Verbatim}{fontsize=\footnotesize,commandchars=\@\^\~,frame=single,label={\textsf{Example}},numbers=left,numbersep=2pt,formatcom=\upshape}
%% A command of the shell
\DefineVerbatimEnvironment{cmdshell}{Verbatim}{commandchars=\@\^\~,fontsize=\footnotesize,formatcom=\upshape}

%%--------------------------------------------------------------------------
%% The reference: the definition of a function, a method, a type, a constant
%% or a macro, its signature on a line and its description indented under it
%%--------------------------------------------------------------------------
\makeatletter
\let\@firsttime=\@undefined
\newlength{\defwidth}
\setlength{\defwidth}{.94\textwidth}
%% A signature: the type, the name in bold, the parameters
\newcommand*{\signature}[3]{%
  \noindent{\ttfamily\raggedright\frenchspacing\hangindent=2em\hangafter=1\relax #1 \textbf{#2}#3\par}}

%% A function: #1 = return type, #2 = name, #3 = parameters
\newcommand{\deffunx}[3]{%
  \defindex{#2}%
  \ifx\@firsttime\@undefined\medskip\fi
  \gdef\@firsttime{}%
  \signature{#1}{#2}{#3}}
\newenvironment{deffun}[3]{%
  \deffunx{#1}{#2}{#3}%
  \begin{flushright}%
  \begin{minipage}{\defwidth}%
}{%
  \end{minipage}%
  \end{flushright}%
  \medskip
  \global\let\@firsttime=\@undefined
}

%% A function template: #1 = template, #2 = return type, #3 = name,
%% #4 = parameters
\newcommand{\defTfunx}[4]{%
  \defindex{#3}%
  \ifx\@firsttime\@undefined\medskip\fi
  \gdef\@firsttime{}%
  \noindent{\ttfamily\raggedright\frenchspacing #1\par}
  \signature{#2}{#3}{#4}}
\newenvironment{defTfun}[4]{%
  \defTfunx{#1}{#2}{#3}{#4}%
  \begin{flushright}%
  \begin{minipage}{\defwidth}%
}{%
  \end{minipage}%
  \end{flushright}%
  \medskip
  \global\let\@firsttime=\@undefined
}

%% A method: #1 = class, #2 = return type, #3 = name, #4 = parameters
\newcommand{\defmethodx}[4]{%
  \defindex{#3}%
  \ifx\@firsttime\@undefined\medskip\fi
  \gdef\@firsttime{}%
  \noindent{\ttfamily\raggedright\frenchspacing\hangindent=2em\hangafter=1\relax\if#2\@empty\else #2 \fi#1::\textbf{#3}#4\par}}
\newenvironment{defmethod}[4]{%
  \defmethodx{#1}{#2}{#3}{#4}%
  \begin{flushright}%
  \begin{minipage}{\defwidth}%
}{%
  \end{minipage}%
  \end{flushright}%
  \medskip
  \global\let\@firsttime=\@undefined
}

%% A static method: #1 = class, #2 = return type, #3 = name, #4 = parameters
\newcommand{\defstaticmethodx}[4]{%
  \defindex{#3}%
  \ifx\@firsttime\@undefined\medskip\fi
  \gdef\@firsttime{}%
  \noindent{\ttfamily\frenchspacing #2 #1::\textbf{#3}#4\hfill\textrm{\textit{static}}}\par}
\newenvironment{defstaticmethod}[4]{%
  \defstaticmethodx{#1}{#2}{#3}{#4}%
  \begin{flushright}%
  \begin{minipage}{\defwidth}%
}{%
  \end{minipage}%
  \end{flushright}%
  \medskip
  \global\let\@firsttime=\@undefined
}

%% A type: #1 = kind (struct, class, enum), #2 = name
\newenvironment{deftype}[2]{%
  \defindex{#2}%
  \medskip
  \noindent{\ttfamily\textbf{#2}\hfill\textrm{\textit{#1}}}\par
  \begin{flushright}%
  \begin{minipage}{\defwidth}%
}{%
  \end{minipage}%
  \end{flushright}%
  \medskip
}

%% A constant: #1 = type, #2 = name
\newenvironment{defconst}[2]{%
  \defindex{#2}%
  \medskip
  \noindent{\ttfamily\frenchspacing #1 \textbf{#2}\hfill\textrm{\textit{const}}}\par
  \begin{flushright}%
  \begin{minipage}{\defwidth}%
}{%
  \end{minipage}%
  \end{flushright}%
  \medskip
}

%% A macro without parameter: #1 = name
\newenvironment{defmacro}[1]{%
  \defindex{#1}%
  \medskip
  \noindent{\ttfamily\textbf{#1}\hfill\textrm{\textit{macro}}}\par
  \begin{flushright}%
  \begin{minipage}{\defwidth}%
}{%
  \end{minipage}%
  \end{flushright}%
  \medskip
}

%% A macro with parameters: #1 = name, #2 = parameters
\newenvironment{defmacroA}[2]{%
  \defindex{#1}%
  \medskip
  \noindent{\ttfamily\textbf{#1}#2\hfill\textrm{\textit{macro}}}\par
  \begin{flushright}%
  \begin{minipage}{\defwidth}%
}{%
  \end{minipage}%
  \end{flushright}%
  \medskip
}

%% The mathematical definition of an operation, one \op per line
\newenvironment{operation}{%
  \vspace*{-5pt}%
  \def\op##1{\noindent{\small\ensuremath{\blacksquare}}\ ##1\par}%
}{%
  \par\smallskip
}
\makeatother

%% The page of the copyright notice
\newenvironment{copyrighttxt}{%
  \parindent0pt%
  \newpage
  \thispagestyle{empty}%
  \null\vfill
}{%
  \newpage
}

%% The numbers of the table of contents and of the lists are wider than in the
%% thesis (3.10.10, 12.13)
\makeatletter
\renewcommand*{\l@section}{\@dottedtocline{1}{1.5em}{2.8em}}
\renewcommand*{\l@subsection}{\@dottedtocline{2}{4.3em}{3.6em}}
\renewcommand*{\l@table}{\@dottedtocline{1}{1.5em}{3.0em}}
\renewcommand*{\l@figure}{\@dottedtocline{1}{1.5em}{3.0em}}
\makeatother

%% A column of a table whose text is set ragged right: L{width}
\newcolumntype{L}[1]{>{\raggedright\arraybackslash}p{#1}}
%% The names of functions and options, in typewriter, are not hyphenated: let
%% the lines holding them stretch rather than overflow into the margin
\setlength{\emergencystretch}{3em}
%% A box that overflows the margin by less than 5 pt, which is 2.5 cm wide, is
%% not worth a warning
\hfuzz=5pt

%%--------------------------------------------------------------------------
%% hyperref comes last
%%--------------------------------------------------------------------------
\usepackage{hyperref}
%% The names of macros and functions, in typewriter, are long and are not
%% hyphenated: a line may break after an underscore (GAOL_PRESERVE_ROUNDING),
%% without a hyphen. In the strings of the PDF the underscore stays itself.
\makeatletter
\let\gaol@underscore\_
\renewcommand*{\_}{\texorpdfstring{\gaol@underscore\hspace{0pt}}{\gaol@underscore}}
\makeatother
\hypersetup{
    unicode      = false,
    pdftoolbar   = true,
    pdfmenubar   = true,
    pdffitwindow = true,
    pdftitle     = {GAOL v5: NOT Just Another Interval Arithmetic Library},
    pdfauthor    = {Jordan Ninin, Frederic Goualard},
    pdfsubject   = {The manual of GAOL v5},
    pdfnewwindow = true,
    pdfkeywords  = {interval arithmetic, IEEE 1788-2015, CORE-MATH,
                    constraint programming, C++},
    colorlinks   = true,
    linkcolor    = blue,
    citecolor    = blue,
    filecolor    = blue,
    urlcolor     = blue,
    breaklinks   = true
}
 loads.
%%
%% The layout is the one of the PhD thesis of Jordan Ninin (Universite de
%% Toulouse, 2010): the class book in 12 points on A4 paper, printed on both
%% sides, the margins of 2.5 cm (top, bottom, outer) and 3.5 cm (inner), the
%% running heads of fancyhdr (the chapter on the even pages, the section on the
%% odd ones, a rule of 1 point under them, the page number between dashes at
%% the foot), the links of hyperref in blue, and the numbered environments
%% Definition, Example, Remark... of the thesis, numbered within the sections.
%% The macros of the reference (\code, \file, \pxref, the environments deffun,
%% defmethod, example, onscreen...) are those of the manual of GAOL 4 by
%% Frederic Goualard (manual/v4/manual.cls), kept under their names so that the
%% text of GAOL 4 the manual takes up, and manual/check_examples.py, which reads
%% the examples, go on working.
%%
%% Copyright (c) 2026 ENSTA, France
%%
%% Created 2026-11-05 by Jordan NININ

%%--------------------------------------------------------------------------
%% Packages
%%--------------------------------------------------------------------------
\usepackage[T1]{fontenc}
\usepackage[utf8]{inputenc}
%% Latin Modern, the outlines of Computer Modern in the T1 encoding: without
%% them pdfTeX embeds bitmaps
\usepackage{lmodern}
\usepackage{textcomp}
\usepackage[english]{babel}
\usepackage{amssymb,amsmath}
\usepackage{graphicx}
\usepackage{xcolor}
\usepackage{colortbl}
\usepackage{multirow}
\usepackage{array}
\usepackage{longtable}
\usepackage{multicol}
\usepackage{fancyvrb}
\usepackage{pifont}
\usepackage[labelsep=endash,labelfont=sc]{caption}
%% The index: \index[idx]{...}, printed by \printindex[idx]
\usepackage{index}
\newindex{idx}{idx}{ind}{Index}

%%--------------------------------------------------------------------------
%% The page
%%--------------------------------------------------------------------------
\raggedbottom
\usepackage[a4paper,top=2.5cm,bottom=2.5cm,left=3.5cm,right=2.5cm,includefoot,includehead,headheight=15pt]{geometry}

\usepackage{fancyhdr}
\pagestyle{fancy}
\fancyhead{}
\fancyfoot{}
\renewcommand{\chaptermark}[1]{\markboth{\chaptername\ \thechapter\ - \ #1}{}}
\renewcommand{\sectionmark}[1]{\markright{\thesection\ - \ #1}}
\renewcommand{\headrulewidth}{1pt}
\fancyhead[LE]{\bfseries\nouppercase{\leftmark}}
\fancyhead[RO]{\bfseries\nouppercase{\rightmark}}
\fancyfoot[CE,CO]{\bfseries--\ \thepage\ --}

%% The first page of a chapter: no running head, the page number at the foot
\fancypagestyle{plain}{%
  \fancyhead{}%
  \fancyfoot[CE,CO]{\bfseries--\ \thepage\ --}%
}
%% A page with nothing
\fancypagestyle{vide}{%
  \fancyhead{}%
  \fancyfoot{}%
  \renewcommand{\headrulewidth}{0pt}%
}
%% The page number alone, under the rule
\fancypagestyle{rien}{%
  \fancyhead{}%
  \fancyfoot[CE,CO]{\bfseries--\ \thepage\ --}%
  \renewcommand{\headrulewidth}{1pt}%
}
%% The running heads of the chapters
\fancypagestyle{YO}{%
  \fancyhead[LE]{\bfseries\nouppercase{\leftmark}}%
  \fancyhead[RO]{\bfseries\nouppercase{\rightmark}}%
  \fancyfoot[CE,CO]{\bfseries--\ \thepage\ --}%
  \renewcommand{\headrulewidth}{1pt}%
}
%% The unnumbered chapters: the same text in both running heads
\newcommand*{\fixedheads}[1]{%
  \fancyhead[LE]{\bfseries\nouppercase{#1}}%
  \fancyhead[RO]{\bfseries\nouppercase{#1}}%
  \fancyfoot[CE,CO]{\bfseries--\ \thepage\ --}%
  \renewcommand{\headrulewidth}{1pt}%
}
\fancypagestyle{INTRO}{\fixedheads{Introduction}}
\fancypagestyle{CONCLU}{\fixedheads{Conclusion}}
\fancypagestyle{BIBLIO}{\fixedheads{Bibliography}}
\fancypagestyle{INDEX}{\fixedheads{Index}}

%% The blank page before a chapter carries nothing
\makeatletter
\def\cleardoublepage{\clearpage
  \if@twoside\ifodd\c@page\else
    \hbox{}\thispagestyle{vide}\newpage\if@twocolumn\hbox{}\newpage\fi\fi\fi}
\makeatother

%% An unnumbered chapter listed in the table of contents
\newcommand{\schapter}[1]{%
  \chapter*{#1}%
  \addcontentsline{toc}{chapter}{#1}%
  \markboth{#1}{#1}%
}

%%--------------------------------------------------------------------------
%% The numbered environments of the thesis: Definition 1.2.1, Example 1.2.2...
%% numbered within the sections, with one counter for all of them
%%--------------------------------------------------------------------------
\newtheorem{Def}{Definition}[section]
\newtheorem{Lem}[Def]{Lemma}
\newtheorem{Prop}[Def]{Proposition}
\newtheorem{Th}[Def]{Theorem}
\newtheorem{Cor}[Def]{Corollary}
\newtheorem{Ex}[Def]{Example}
\newtheorem{Rem}[Def]{Remark}
\def\demo{{\em Proof. }}
\def\boite{\hfill{$\hbox{\vrule height4pt width4pt}$}}

%%--------------------------------------------------------------------------
%% Mathematics
%%--------------------------------------------------------------------------
\newcommand{\RSet}{\ensuremath{\mathbb{R}}}
\newcommand{\FSet}{\ensuremath{\mathbb{F}}}
\newcommand{\ISet}{\ensuremath{\mathbb{I}}}
\newcommand{\itv}[2]{\ensuremath{[#1,\,#2]}}
\newcommand{\roundDn}[1]{\ensuremath{\downarrow#1\downarrow}}
\newcommand{\roundUp}[1]{\ensuremath{\uparrow#1\uparrow}}
\newcommand{\roundNearest}[1]{\ensuremath{\updownarrow#1\updownarrow}}
\newcommand{\hull}[1]{\ensuremath{\mathop{\Box} #1}}
\newcommand{\leftBound}[1]{\ensuremath{\underline{#1}}}
\newcommand{\rightBound}[1]{\ensuremath{\overline{#1}}}
\newcommand*{\prevdouble}[1]{\ensuremath{{#1}^-}}
\newcommand*{\nextdouble}[1]{\ensuremath{{#1}^+}}
\newcommand{\returns}{\hbox{\ding{229}}}
\let\leq=\leqslant
\let\geq=\geqslant
\let\emptyset=\varnothing

%%--------------------------------------------------------------------------
%% Names in the text
%%--------------------------------------------------------------------------
%% Code, a name of the program
\newcommand*{\code}[1]{\textup{\texttt{#1}}}
%% A name to stand out: a value of an enumeration, a format
\newcommand*{\element}[1]{\textup{\textbf{\texttt{#1}}}}
%% A definition, in the text
\newcommand{\dfn}[1]{\textsl{#1}}
%% The name of a command or a program
\newcommand{\cmd}[1]{\textup{\texttt{#1}}}
%% The name of an option, and of a long option (-{}-name)
\newcommand{\option}[1]{\textup{\texttt{#1}}}
\newcommand{\loption}[1]{\textup{\texttt{-{}-#1}}}
%% The name of a variable or a parameter
\newcommand{\var}[1]{\textsl{#1}}
%% The name of a file
\newcommand{\file}[1]{\textup{\texttt{#1}}}
%% The name of an environment variable
\newcommand{\env}[1]{\textup{\texttt{#1}}\index[idx]{#1@\texttt{#1}}}
%% A name of IEEE 1788-2015
\newcommand*{\ieee}[1]{\textsf{#1}}
%% A non terminal and a terminal of a grammar
\newcommand*{\NT}[1]{\uppercase{\textit{#1}}}
\newcommand*{\TM}[1]{\textsf{#1}}
%% What GAOL v5 changes: a mark in the text (process.md, section 4)
\newcommand*{\newinvfive}{\unskip\nobreak\ \textcolor{blue}{\textsf{\scriptsize[GAOL~v5]}}}
%% A comment in a program, and what a program prints
\definecolor{remark}{gray}{0.40}
\newcommand{\rem}[1]{\textcolor{remark}{\textsl{#1}}}
\newcommand{\outputs}[1]{\textcolor{remark}{\textsl{\returns{} #1}}}
%% A warning, in the text: the manual writes its warnings as Remarks, this
%% command stays for the sentences that need a word
\newcommand*{\important}{\textbf{Caution.}~}

%%--------------------------------------------------------------------------
%% Cross references and citations
%%--------------------------------------------------------------------------
\newcommand*{\pxref}[1]{see Section~\hyperref[#1]{\ref{#1}}, page~\pageref{#1}}
\newcommand*{\spxref}[1]{see Section~\hyperref[#1]{[\nameref{#1}]}, page~\pageref{#1}}
\newcommand*{\xref}[1]{See Section~\hyperref[#1]{\ref{#1}}, page~\pageref{#1}}
\newcommand*{\sxref}[1]{See Section~\hyperref[#1]{[\nameref{#1}]}, page~\pageref{#1}}
\newcommand*{\exref}[1]{Section~\hyperref[#1]{\ref{#1}}, page~\pageref{#1}}
\newcommand*{\sexref}[1]{Section~\hyperref[#1]{[\nameref{#1}]}, page~\pageref{#1}}
%% The citations of the manual of GAOL 4 were printed in the margin; the
%% manual of GAOL v5 cites in the text, as the thesis does
\makeatletter
\newcommand*{\margincite}{\@ifstar{\margincite@}{\margincite@}}
\newcommand*{\margincite@}[2][]{%
  \def\@tempa{#1}%
  \ifx\@tempa\@empty\cite{#2}\else\cite[#1]{#2}\fi}
\newcommand*{\marginnocite}[1]{\nocite{#1}}
\makeatother

%%--------------------------------------------------------------------------
%% The index
%%--------------------------------------------------------------------------
\newcommand*{\findex}[1]{\index[idx]{#1@\texttt{#1}}}
\newcommand*{\cindex}[1]{\index[idx]{#1}}
\newcommand*{\idxself}{\~{}}
\newcommand*{\hyperpagebf}[1]{\textbf{\hyperpage{#1}}}
%% The main entry of a name, whose page number is in bold
\newcommand*{\defindex}[1]{\index[idx]{#1@\texttt{#1}|textbf}}

%%--------------------------------------------------------------------------
%% Programs, commands and outputs (fancyvrb): @name^argument~ is the command
%% \name{argument}, and manual/check_examples.py reads the lines @outputs^...~
%% of the examples
%%--------------------------------------------------------------------------
%% What a program prints
\DefineVerbatimEnvironment{onscreen}{Verbatim}{commandchars=\@\^\~,frame=single,label={\textsf{Output}},fontsize=\footnotesize,formatcom=\upshape}
%% An example program
\def\example{\medskip\examplevb}
\def\endexample{\endexamplevb}
\DefineVerbatimEnvironment{examplevb}{Verbatim}{fontsize=\footnotesize,commandchars=\@\^\~,frame=single,label={\textsf{Example}},numbers=left,numbersep=2pt,formatcom=\upshape}
%% A command of the shell
\DefineVerbatimEnvironment{cmdshell}{Verbatim}{commandchars=\@\^\~,fontsize=\footnotesize,formatcom=\upshape}

%%--------------------------------------------------------------------------
%% The reference: the definition of a function, a method, a type, a constant
%% or a macro, its signature on a line and its description indented under it
%%--------------------------------------------------------------------------
\makeatletter
\let\@firsttime=\@undefined
\newlength{\defwidth}
\setlength{\defwidth}{.94\textwidth}
%% A signature: the type, the name in bold, the parameters
\newcommand*{\signature}[3]{%
  \noindent{\ttfamily\raggedright\frenchspacing\hangindent=2em\hangafter=1\relax #1 \textbf{#2}#3\par}}

%% A function: #1 = return type, #2 = name, #3 = parameters
\newcommand{\deffunx}[3]{%
  \defindex{#2}%
  \ifx\@firsttime\@undefined\medskip\fi
  \gdef\@firsttime{}%
  \signature{#1}{#2}{#3}}
\newenvironment{deffun}[3]{%
  \deffunx{#1}{#2}{#3}%
  \begin{flushright}%
  \begin{minipage}{\defwidth}%
}{%
  \end{minipage}%
  \end{flushright}%
  \medskip
  \global\let\@firsttime=\@undefined
}

%% A function template: #1 = template, #2 = return type, #3 = name,
%% #4 = parameters
\newcommand{\defTfunx}[4]{%
  \defindex{#3}%
  \ifx\@firsttime\@undefined\medskip\fi
  \gdef\@firsttime{}%
  \noindent{\ttfamily\raggedright\frenchspacing #1\par}
  \signature{#2}{#3}{#4}}
\newenvironment{defTfun}[4]{%
  \defTfunx{#1}{#2}{#3}{#4}%
  \begin{flushright}%
  \begin{minipage}{\defwidth}%
}{%
  \end{minipage}%
  \end{flushright}%
  \medskip
  \global\let\@firsttime=\@undefined
}

%% A method: #1 = class, #2 = return type, #3 = name, #4 = parameters
\newcommand{\defmethodx}[4]{%
  \defindex{#3}%
  \ifx\@firsttime\@undefined\medskip\fi
  \gdef\@firsttime{}%
  \noindent{\ttfamily\raggedright\frenchspacing\hangindent=2em\hangafter=1\relax\if#2\@empty\else #2 \fi#1::\textbf{#3}#4\par}}
\newenvironment{defmethod}[4]{%
  \defmethodx{#1}{#2}{#3}{#4}%
  \begin{flushright}%
  \begin{minipage}{\defwidth}%
}{%
  \end{minipage}%
  \end{flushright}%
  \medskip
  \global\let\@firsttime=\@undefined
}

%% A static method: #1 = class, #2 = return type, #3 = name, #4 = parameters
\newcommand{\defstaticmethodx}[4]{%
  \defindex{#3}%
  \ifx\@firsttime\@undefined\medskip\fi
  \gdef\@firsttime{}%
  \noindent{\ttfamily\frenchspacing #2 #1::\textbf{#3}#4\hfill\textrm{\textit{static}}}\par}
\newenvironment{defstaticmethod}[4]{%
  \defstaticmethodx{#1}{#2}{#3}{#4}%
  \begin{flushright}%
  \begin{minipage}{\defwidth}%
}{%
  \end{minipage}%
  \end{flushright}%
  \medskip
  \global\let\@firsttime=\@undefined
}

%% A type: #1 = kind (struct, class, enum), #2 = name
\newenvironment{deftype}[2]{%
  \defindex{#2}%
  \medskip
  \noindent{\ttfamily\textbf{#2}\hfill\textrm{\textit{#1}}}\par
  \begin{flushright}%
  \begin{minipage}{\defwidth}%
}{%
  \end{minipage}%
  \end{flushright}%
  \medskip
}

%% A constant: #1 = type, #2 = name
\newenvironment{defconst}[2]{%
  \defindex{#2}%
  \medskip
  \noindent{\ttfamily\frenchspacing #1 \textbf{#2}\hfill\textrm{\textit{const}}}\par
  \begin{flushright}%
  \begin{minipage}{\defwidth}%
}{%
  \end{minipage}%
  \end{flushright}%
  \medskip
}

%% A macro without parameter: #1 = name
\newenvironment{defmacro}[1]{%
  \defindex{#1}%
  \medskip
  \noindent{\ttfamily\textbf{#1}\hfill\textrm{\textit{macro}}}\par
  \begin{flushright}%
  \begin{minipage}{\defwidth}%
}{%
  \end{minipage}%
  \end{flushright}%
  \medskip
}

%% A macro with parameters: #1 = name, #2 = parameters
\newenvironment{defmacroA}[2]{%
  \defindex{#1}%
  \medskip
  \noindent{\ttfamily\textbf{#1}#2\hfill\textrm{\textit{macro}}}\par
  \begin{flushright}%
  \begin{minipage}{\defwidth}%
}{%
  \end{minipage}%
  \end{flushright}%
  \medskip
}

%% The mathematical definition of an operation, one \op per line
\newenvironment{operation}{%
  \vspace*{-5pt}%
  \def\op##1{\noindent{\small\ensuremath{\blacksquare}}\ ##1\par}%
}{%
  \par\smallskip
}
\makeatother

%% The page of the copyright notice
\newenvironment{copyrighttxt}{%
  \parindent0pt%
  \newpage
  \thispagestyle{empty}%
  \null\vfill
}{%
  \newpage
}

%% The numbers of the table of contents and of the lists are wider than in the
%% thesis (3.10.10, 12.13)
\makeatletter
\renewcommand*{\l@section}{\@dottedtocline{1}{1.5em}{2.8em}}
\renewcommand*{\l@subsection}{\@dottedtocline{2}{4.3em}{3.6em}}
\renewcommand*{\l@table}{\@dottedtocline{1}{1.5em}{3.0em}}
\renewcommand*{\l@figure}{\@dottedtocline{1}{1.5em}{3.0em}}
\makeatother

%% A column of a table whose text is set ragged right: L{width}
\newcolumntype{L}[1]{>{\raggedright\arraybackslash}p{#1}}
%% The names of functions and options, in typewriter, are not hyphenated: let
%% the lines holding them stretch rather than overflow into the margin
\setlength{\emergencystretch}{3em}
%% A box that overflows the margin by less than 5 pt, which is 2.5 cm wide, is
%% not worth a warning
\hfuzz=5pt

%%--------------------------------------------------------------------------
%% hyperref comes last
%%--------------------------------------------------------------------------
\usepackage{hyperref}
%% The names of macros and functions, in typewriter, are long and are not
%% hyphenated: a line may break after an underscore (GAOL_PRESERVE_ROUNDING),
%% without a hyphen. In the strings of the PDF the underscore stays itself.
\makeatletter
\let\gaol@underscore\_
\renewcommand*{\_}{\texorpdfstring{\gaol@underscore\hspace{0pt}}{\gaol@underscore}}
\makeatother
\hypersetup{
    unicode      = false,
    pdftoolbar   = true,
    pdfmenubar   = true,
    pdffitwindow = true,
    pdftitle     = {GAOL v5: NOT Just Another Interval Arithmetic Library},
    pdfauthor    = {Jordan Ninin, Frederic Goualard},
    pdfsubject   = {The manual of GAOL v5},
    pdfnewwindow = true,
    pdfkeywords  = {interval arithmetic, IEEE 1788-2015, CORE-MATH,
                    constraint programming, C++},
    colorlinks   = true,
    linkcolor    = blue,
    citecolor    = blue,
    filecolor    = blue,
    urlcolor     = blue,
    breaklinks   = true
}
